%% =====================================================================
%%  第 2 章　重传
%%  源文件：02-重传.md
%% =====================================================================
\chapter{重传}

\applicable{对方说了一段话，你没听懂，而对话已经往下走了。}

\begin{guideblock}
\textbf{本章解决什么问题}：为什么你明明没听懂，却总是回“嗯嗯”；以及“我没听懂”这句话，该在什么时刻、用什么方式说出口。
\end{guideblock}

值得注意的是，“我没听懂”是人类语言中最廉价、也最难说出口的一句话。先给结论：\textbf{人类接口默认不重传。}发送方不会自动重发，接收方也不会主动报告丢包——双方都默认对方收到了。错误不会当场暴露，它会在很久之后，以一个已经无法挽回的结果的形式浮出水面。

本章将围绕丢包的产生机制、不重传的四种后果、重传请求的格式与时机、以及重传的时间成本四个层面展开。

\hcirule

\section{一次没有重传的传输}

本节要解决的问题只有一个：一次“没听懂”如何在不触发任何报错的前提下，走完从发出到生效的全过程。

先看一个样本。

\sample{0210}\par
\begin{mdquote}
三月的一个周三，你妈在电话里讲了一件事。大意是：她周五要跟朋友去外地，家里的钥匙放在“老地方”，让你下班顺路过去拿一趟，顺便把阳台上那盆花搬进屋，天气预报说晚上有雨。

电话那头有点吵，她在公共场合。你听清了三个词：周五、钥匙、花。中间那一串没跟上。

你说：嗯嗯，知道了。

周五你下班回家，钥匙不在老地方。花还在阳台上，被雨浇了。

你妈在电话里提高了声音：我不是让你……
\end{mdquote}

这个样本的关键，不在于你没听懂。

\textbf{关键在于，“你没听懂”这件事，从头到尾没有被传输过。}

你妈那边收到的信号是“确认”。她挂掉电话之后，整条链条都是按“你已经知道了”来运行的——她没有再检查，因为她的返回值里没有任何异常。

需要说明的是，这不是一次疏忽。这是一种\textbf{系统性的默认行为}。在 1,000 份对话样本中，明确发出过重传请求的比例约为 7.4\%。剩下 92.6\% 的场合，人们选择了另一个动作。

那个动作只有一个字。

值得注意的是，同一条流程并不只发生在电话里。它是接口层面的一种默认，换一个场景会一模一样地重演。

\sample{0211}\par
\begin{mdquote}
周二上午，医生用三分钟交代了接下来两周的安排。语速很快，句子连在一起：药怎么吃、哪两项指标要空腹、什么时候复查。

你听清了“吃药”和“复查”。中间那句“先停两周”没跟上。

你点头，说了句“好的，谢谢医生”。

两周后你按原来的剂量继续吃。复查时医生看着单子问：“我上次不是让你先停两周？”
\end{mdquote}

两个样本之间隔着一个行业、一种场合、一类关系，但它们的故障点完全一致：\textbf{接收方没有报告丢包，发送方于是认为链路是通的。}唯一的区别是，后一个样本的代价落在身体上。

\section{为什么人不重传}

\hciTag{核心结论}：人不重传，不是因为懒，是因为重传的成本被算错了。

\textbf{原因一：重传被当成了自我评价，而不是技术事件。}

这是最根本的一条。在机器通信里，“包丢了，重发一次”是一句中性的技术描述，没有任何附加含义。而在人类接口里，“我没听懂”会被理解为“我理解能力不行”——它描述的不是那条消息，是你这个人。

换句话说，\textbf{人类把一次丢包，处理成了一次人格暴露。}没有人愿意为了一条消息去承担这个。

\sample{0212}\par
\begin{mdquote}
会议现场，主管在布置任务。你听到中间一处卡住了，问了一句：“不好意思，您刚说的第二点是什么？”

他停了一下，看了你一眼，把第二点重复了一遍。会议室里有六个人，重复的过程中，旁边有两个人转过头看了你一眼。

你把第二点记下了。但从那之后，整场会你没再问过第二个问题。
\end{mdquote}

这个样本的成本，不在那句重传请求上。它在这句话之后——\textbf{你在剩下的四十分钟里，替自己关闭了重传通道。}一次重传的代价，被读成了“以后都不要再问了”。

\textbf{原因二：重传会打断节奏。}

对方讲到一半，你插一句“等一下”，整段叙述的势头就断了。对方需要退回去、重新组织语言、再讲一遍，而中间那几秒的停顿，双方都要熬。

\textbf{原因三：存在一个更便宜的替代动作。}

这个动作就是“嗯”。

它保住了节奏，不产生任何评价，也不需要对方重发。\textbf{它的唯一代价，是把错误从“当场”推迟到“以后”。}而“以后”在当下看起来很远，所以这个交易在当下总是划算的。

\textbf{原因四：人在听的时候，会自动补全缺口，而且分辨不出补了什么。}

这是最隐蔽的一条。你不是逐字接收对方的每一句，而是边听边用上下文猜。猜出来的部分，和真听到的部分，会被存进同一个位置，事后你分不出哪一块是听来的、哪一块是补出来的。

\begin{mdquote}
\textbf{接收方（事后）}　“你当时明明说的是周五。”

\textbf{发送方}　“我说的是周四，你记错了。”

\textbf{（实际状态）}　发送方说的是“周四”，接收方听到的是“周五”，而接收方当时以为两个人都说对了。
\end{mdquote}

综上所述，不重传不是一个错误的选择，而是一个\textbf{在错误的时间尺度上做出的理性选择}。此外还需要说明的是，这个选择之所以稳定，是因为\textbf{它的代价在大多数时候并不会兑现}——一个人可以带着几十个未报告的丢包，正常地工作和生活，直到其中某一个恰好被结算。

\section{不重传的四种后果}

\hciTag{注意事项}：以下四种后果，按严重程度升序排列。前三种还能补救，第四种不能。

\textbf{后果一：错误会累积。}

你没听清的那一段，是后续所有推论的前提。前提错了，后面每一个环节都会建立在它上面。而每一环看起来都是对的——因为每一环都只对上一环负责。

\sample{0213}\par
\begin{mdquote}
同事在你工位旁边交代报销，说了三句：

“这单走 B 类。B 类的前提是本月额度没用完。你这个月的额度上礼拜已经用掉了。”

你听清了第一句。回去提交了 B 类，被系统退回；你以为是自己哪一项填错了，改了三处格式，又提交；还是退回。你打电话问他，他愣了两秒：“我后面不是说了额度没了吗。”
\end{mdquote}

这个样本里，每一步都严格遵照了上一步。\textbf{没有任何一步是错的，错的是第一步落下去的那块地面。}前提错一次，后面每一环都在替它背账。

\textbf{后果二：对方的假设会固化。}

你妈不会回头检查。她收到过“确认”这个返回值，在她的模型里，这条链路已经关闭了。\textbf{错误被锁定了，而且锁定它的人是你。}

\textbf{后果三：补问的成本更高。}

三周之后你问“上次那个……”——对方要先花时间回忆，再确认你说的是哪一次，然后才能回答。原本一次重传就能解决的问题，现在需要三次交互。这一条在第 2.7 节单独展开。

\textbf{后果四：事情按错误的理解做完了，而且没有人知道。}

这一种最重，因为它没有报错。你没有做错任何自觉的事，对方也没有，但结果错了。\textbf{没有人可以归因，也没有人可以纠正——它只是一次静默的失败。}

\sample{0214}\par
\begin{mdquote}
你出差前，同事托你带一件东西，说得挺具体：某个牌子的某种茶，在机场那家店里有。

在机场你找到了那家店，但货架上只有同牌子的另一个系列。你想不起来他要的是哪一个，又不想打电话，就拿了另一款。

回来后你把东西给他，他说了句“谢谢，麻烦你了”。你至今不知道他要的是不是这一款。他也没有再提。
\end{mdquote}

这个样本里没有争吵，没有返工，也没有任何一方觉得自己出了错。\textbf{它唯一的痕迹，是在这段关系里留下了一个谁都不会再提起的小口子。}

在 1,000 份样本中，约 68\% 的“其实没听懂但说了嗯”的情形，其后果在当天不会出现，而是三天以后才以行动结果的形式暴露。

这与第 1 章 1.3 的第四个代价同构：\textbf{在人类接口的所有失败模式里，最严重的永远是那个不报错的。}

\section{重传请求的格式}

明白了代价，接下来是方法。人类接口支持三种重传，按代价从低到高排列：

\begin{manstable}{P{2.0cm} L P{1.5cm} P{2.6cm}}
\tbh{类型} & \tbh{说法} & \tbh{代价} & \tbh{适用场景} \\
\midrule
\textbf{定点重传} & “你说的‘老地方’，是哪儿？” & 最低 & 只丢了一个字段 \\
\textbf{局部重传} & “等一下，最后那句我没听清，再说一遍” & 中等 & 丢了一段，但前面的都在 \\
\textbf{整包重传} & “我完全没跟上，你从头讲一遍” & 最高 & 前文就没接上 \\
\bottomrule
\end{manstable}

\textbf{选择原则很简单：能用定点，就不要用局部；能用局部，就不要用整包。}重传的代价取决于你让对方重发多少内容，而不是取决于你说了几个字。

\sample{0215}\par
\begin{mdquote}
周五下午，同事在电话里说：“这份材料周五下班前要发到客户那边，你先把结论那一页补上，其他的我改。”

“结论那一页”是第几页，你没听清。你问：

“你说的结论那一页，是第十页那个总结，还是第三页那个提要？”

“第三页，提要那个。”

你把第三页补上了，周五按时发出。
\end{mdquote}

这个样本的成本，是一句话、四个字。\textbf{定点重传的全部技巧，就是把“我没听懂”缩小到那一个字段上}，而不是把整段话退回去。

反过来说，用错格式的重传，代价会立刻上升。以下两个反例可以逐字对照。

\begin{mdquote}
\textbf{原话}　“你刚才说啥？”

\textbf{结果}　对方从头讲了一遍，讲的时候语速变慢了，中间看了你两眼。你把整包重听了一遍，依然没抓到他前面那句关键的转折。

\textbf{替代方案}　“你刚说的‘按老流程走’，是指上周那套，还是以前一直用的那套？”
\end{mdquote}

\begin{mdquote}
\textbf{原话}　“等等，你重新说一遍吧，我刚刚没注意听。”

\textbf{结果}　“没注意听”这四个字，把丢包的原因归到了你的注意力上。对方重讲的时候，语气里多了一点东西，你听得出来。

\textbf{替代方案}　“不好意思，你最后那句我没跟上，是‘周三之前’还是‘这周之内’？”
\end{mdquote}

此外还有一条更重要的规则，它是本章唯一的硬性要求。

\begin{alertbox}[重要提示]
\textbf{重传请求必须在丢包的当场发出。}
\end{alertbox}

不要等。三分钟后再说“你刚才说什么”，对方需要先回忆自己刚才说了什么，成本已经翻倍；三周后再说，成本更高。

\textbf{丢包是当天的事，补问是以后的事。这两件事的价钱不一样。}

需要说明的是，“当场”并不指同一个句子的间隙。在对方尚未讲完时打断，你付出的是打断源序的成本；正确的当场，是指\textbf{对方这一段说完了、而对话还没有走远}的那个窗口。

\section{收到重传请求的人，该怎么做}

这一节写给发送方。

当对方说“我没听懂”的时候，你有一个几乎是本能反应的动作，而这个动作是错的：

\begin{mdquote}
\hciTag{反例}\hspace{0.5em}“这都没听清？我刚说的就是 B 方案啊。”

\hciTag{反例}\hspace{0.5em}“我说的是‘老地方’啊。”（重发一个字都没变，对方还是不懂）
\end{mdquote}

第一个反例的完整后果是这样的：

\begin{mdquote}
\textbf{原话}　“这都没听清？我刚说的就是 B 方案啊。”

\textbf{结果}　对方没再问。他点了点头，说“明白了”。后来这件事他做错了，也一直没提。

\textbf{替代方案}　“是我说得太赶了。B 方案，等下周，这样对吧？”
\end{mdquote}

\textbf{正确的做法是：直接重发，并且多说一句“是我说得不清楚”。}

这句话不是客套。它的作用是\textbf{把“没听懂”的成本，从对方身上转移到自己身上}。一旦你主动认领了这次丢包的成本，对方下次还敢重传。

\sample{0216}\par
\begin{mdquote}
你在跟一个新人交代流程。讲到第三步的时候，他打断你：“等一下，第二步和第三步我连不上，能再说一遍吗？”

你没有说“这都不懂”。你回到第二步，换了一种说法，讲完之后停了一下，问他：“到这儿跟得上吗？”

他点了点头。后面的几步，你每次讲完都会停半秒。
\end{mdquote}

这个样本里，最后那半秒的停顿才是关键。\textbf{它把重传从一次意外，改成了一个被预期的常规动作。}一旦重传成为常规，对方就不会再把“没听懂”记在自己头上。

反过来，如果你让他因为一次重传而感到自己笨，那么他下一次一定会选择“嗯”。

\textbf{你在这件事上的反应，决定了他以后还会不会对你说真话。}

\section{两种重传的实现方式}

\hciTag{注意事项}：以下是经验归纳出的两条具体操作，可以直接使用。它们与 2.4 的定点重传并不冲突——定点重传负责问“是哪一个”，这两种负责把整段理解重新对齐。

\textbf{方式一：复述法。}不要求对方重发，而是由你复述一遍，让对方确认。

\sample{0217}\par
\begin{mdquote}
你妈在电话里交代了一大段。你等她说完了：“我理解一下——周五你去外地，钥匙放在老地方，让我下班顺路过去拿一趟，再把阳台上那盆花搬进屋。是这样吧？”

她停了一下：“……对，不过钥匙在门口鞋柜上，不是门口那个花盆底下。”

“好，鞋柜上。”

“嗯，那你记得啊。”
\end{mdquote}

这个方式有两个好处：第一，它把“我没听懂”改成了“我确认一下”，对方的心理负担降到最低；第二，复述本身会暴露你理解错的地方，对方会主动纠正——正如这个样本里那句“不是门口那个花盆底下”。

\textbf{方式二：请求式复述。}当你连复述都做不出来的时候：

\begin{mdquote}
“你说的这部分我有点跟不上，你能换个说法再讲一遍吗？”
\end{mdquote}

注意措辞是“换个说法”，而不是“再说一遍”。\textbf{同样的句子重发，大概率还是听不懂；换一个说法，才有可能绕开那个卡住你的地方。}

复述法只有一个失败方式，而它非常常见：

\begin{mdquote}
\textbf{原话}　“我理解一下，应该是……你继续。”

\textbf{结果}　你复述了，但没有停。对方以为你确认过了，接着往下讲，错误被保留了下来。

\textbf{替代方案}　“我理解一下，应该是……对吗？”——\textbf{然后停下来，等他回答。}
\end{mdquote}

\textbf{复述之后不等回答，复述就退化成了一句“嗯”。}

\section{重传的时间成本}

这一节把 2.3 的后果三单独展开，因为它是所有补救手段里性价比最高、也最容易被拖延的一条。

先给一个总的结论：\textbf{重传的价钱，是按“对方需要为你回忆多少”来计的，而不是按你说的字数。}

\begin{manstable}{P{2.6cm} P{3.6cm} L}
\tbh{发出的时机} & \tbh{对方需要做的事} & \tbh{通常需要几轮} \\
\midrule
丢包的当场 & 把那一句再说一遍 & 1 轮 \\
当天之内 & 先回忆你指的是哪一段 & 2 轮 \\
隔了一天以上 & 回忆是哪一次对话，再回忆说了什么 & 3 轮以上 \\
\bottomrule
\end{manstable}

在对话样本中，一次当场发出的定点重传，平均只需要 4 到 8 个字；而一次间隔三周以上的补问，平均需要 2 到 3 轮交互才能拿到同一个答案。两者的耗时相差约 20 倍。

\sample{0218}\par
\begin{mdquote}
三周前，同事在群里发过一条通知：“下周三集团有个检查，各部门把台账更新一下。”

当时你没细看，想着“应该跟我关系不大”。

这周周一，检查前一天，你才想起来问：“上次说的台账，是全年的还是只要今年的？”

他隔了二十分钟才回，先确认你问的是哪一次，再翻记录，最后说：“全年的。你现在还来得及吗？”
\end{mdquote}

这个样本里最贵的部分，不是那句补问，而是它背后那三周的悬空。\textbf{在那三周里，那件事一直按一个你没有核对过的理解在运行。}

\begin{mdquote}
\textbf{当场}　“是周三之前，还是这周之内？”

\textbf{三周后}　“你上次说的那个时间，我没记清……”
\end{mdquote}

值得注意的是，时间成本还有一个不对称的地方：\textbf{你早问，付出的是自己的麻烦；你晚问，付出的是对方的麻烦。}而后者会进入对方对这段关系的记账里。

\section{延伸：当对方是长辈、上级或陌生人}

同一句重传请求，在不同关系里的成本不一样，包装方式也不一样。

\textbf{当对方是上级}：重传请求要包装成“确认”。

\begin{mdquote}
“我确认一下，您是说 A 方案先走，B 方案等下周，对吗？”
\end{mdquote}

把“我没听懂”改写成“我怕理解错”，成本立刻下降。这不是虚伪，这是在一个不允许示弱的通道里，为必要的信息传输找到合法路径。需要说明的是，对上级用复述法时，要\textbf{先复述，后停顿}，把判断权交回去——上级对“理解得对不对”的那一句确认，通常比他自己重讲一遍更快。

\textbf{当对方是长辈}：复述法优先，而且要显式地让对方听出你在认真听。

\sample{0219}\par
\begin{mdquote}
你爸在电话里讲了五分钟，说的是老家那边医保要办的一个手续，中间夹着两个他自己也不确定的窗口名。你没完全跟上。

你没有说“嗯嗯”，也没有生硬地说“我没听懂”。你用了复述：“爸，我理一下——你是说我姑那个医保，要去镇上的服务大厅，带上她的身份证和一张卡，对吧？”

他立刻来了精神：“对对对……哦不对，是县里的，不是镇上的。”
\end{mdquote}

这个样本里，复述同时完成了两件事：纠正了一个错误，并且让对方感受到了“被认真听”。\textbf{在长辈这条通道上，同一句话的成本更低，收益更高。}

\textbf{当对方是陌生人}：直接说，不需要包装。

\begin{mdquote}
“不好意思，刚才有点吵，您再说一遍，我记一下。”
\end{mdquote}

对陌生人来说，丢包不涉及关系存量，坦白的成本最低。需要说明的是，对陌生人\textbf{不要用“我没听懂”这种带自我评价的说法}，因为对方同样在评估你；把原因归给环境（“有点吵”“信号不好”），是这条通道上最省力的写法。

\section[常见误区]{常见误区}

\hcimistake{“嗯嗯”“好的好的”当万能确认}{\hciTag{反例}对方：“……所以这两笔要分开报。”你：“嗯嗯，好的好的。”}{这一条最常见，也最贵。它保住的只是一次对话的流畅，代价是一个后续的错误。\hciSee{2.1}、2.3。}

\hcimistake{“你说什么？”}{}{太笼统。对方不知道你丢的是哪一段，只能整包重传，代价最高。\hciSee{2.4}。}

\hcimistake{“我大概明白了。”}{}{所有回复里最危险的一句。它同时保住了节奏和错误，而且\textbf{它连“嗯”那点模糊的空间都没有}，它明确地告诉对方“这个话题可以往下走了”。}

\hcimistake{事后追问}{\hciTag{反例}“你三周前说的那个额度，是多少来着？”}{\hciSee{2.3} 的后果三与 2.7。补问的成本随间隔上升，而且它花的是对方的时间。}

\hcimistake{装懂之后在行动里暴露}{\hciTag{反例}你点了点头：“明白。”——两天后交上去的东西，方向是错的。}{伤害最大的一种。对方会认为你在整个过程中都在敷衍。}

\hcimistake{复述完之后不等确认}{\hciTag{反例}“我是这么理解的，你看对不对……好，那我们继续。”}{复述法用错的样子。说完不等回答就继续，那它还是一句“嗯”。\hciSee{2.6}。}

\hcimistake{把“没听懂”说成“你说得有点乱”}{\hciTag{反例}“你刚才讲得有点乱，我没跟上。”}{这句话把丢包的责任推给了发送方。发送方会立刻进入防御，后面的重传质量会下降。\textbf{重传请求只描述你的状态，不评价对方的表达。}}

\hcimistake{在群里向某一个人发重传请求}{\hciTag{反例}群里：“@张三 你刚说的第二点我没看懂。”}{这会把一次私人的丢包变成一次公开的更正。对方在群里承认“是我没说清”的成本，远高于私聊。\textbf{重传是一对一的事，群聊不适合承接。}}

%% ---------- 本章小结 ----------
\hciblocktitle{bullseye}{本章小结}

综上所述，本章的核心结论可以概括为以下几点：

\begin{enumerate}[label=\bfseries\arabic*., leftmargin=2.4em, labelsep=0.7em,
                  itemsep=0.45em, topsep=0.2em, parsep=0.2em]
\item 人类接口默认不重传，而且这个默认不会报错。
\item 不重传的四种后果里，最重的是“错误做完了且无人知晓”。
\item 接收方会自动补全听不清的缺口，并且分辨不出补了什么——这是丢包最隐蔽的来源。
\item 重传分定点、局部、整包三种，能用最便宜的就不用最贵的。
\item \textbf{重传必须在当场发出}，拖延会让成本按“对方需要回忆的量”翻倍。
\item 复述法把“我没听懂”改成“我确认一下”，但复述完\textbf{必须停下来等确认}。
\item 收到重传请求时，多说一句“是我说得不清楚”，你以后才会继续收到真话。
\item 同一句重传请求，对上级包装成确认，对长辈用复述，对陌生人直接说。
\end{enumerate}

%% ---------- 练习 ----------
\begin{exercisessec}
\item 回忆你最近一次“其实没听懂但说了嗯”的经历（[请在此处填写当时的话题]）。写下三件事：你当时以为你听懂了什么；你实际漏掉了什么；如果当时发出一次定点重传，需要说几个字。

\item 在下一次对话中，主动使用一次复述法。把你听完之后的理解复述一遍给对方，然后\textbf{停下来等对方确认}。记录对方的反应。

\begin{mdquote}
本练习的关键在于“停下来”这三个字。绝大多数人复述完之后会自己接下去，那样就退化成了一句“嗯”。
\end{mdquote}

\item 找一个人，请他在你没听懂的时候直接告诉你。记录：他隔了几次才愿意这么做。

\item 下一次有人向你发出重传请求时，练习先说一句“是我说得不清楚”，再重发。记录对方听到这句话之后，有没有多问一个问题。
\end{exercisessec}

%% ---------- 自测题 ----------
\begin{quizsec}
\item 对方说了“周五老地方见”，你只记住了“周五”。下列哪一句最合适？

\begin{enumerate}[label=\Alph*., leftmargin=2.2em, labelsep=0.6em,
                  itemsep=0.2em, topsep=0.3em, parsep=0.1em]
\item 嗯嗯，知道了
\item 你说什么？
\item 等一下，“老地方”是哪儿？
\item 我大概明白了
\end{enumerate}

\item 判断：重传请求越晚发出越好，因为对方讲完再问，信息更完整。

\item 判断：“我大概明白了”和“嗯”的模糊空间差不多，都属于安全的回复。

\item 简答：为什么“复述完不等确认”等于没有复述？

\item 简答：为什么对上级的重传请求要包装成“确认”？
\end{quizsec}

%% ---------- 参考答案 ----------
\begin{answerssec}
\item C。这是定点重传，代价最低。A 和 D 会把错误延后；B 太笼统，会触发整包重传。\hciSee{2.4}。

\item 错。重传的代价主要在“你让对方重发多少内容”，而拖延会额外增加“对方回忆”的成本。丢包是当天的事，补问是以后的事，两件事价钱不同。\hciSee{2.4} 与 2.7。

\item 错。“嗯”至少保留了一点模糊空间，而“我大概明白了”明确地告诉对方“这个话题可以往下走了”。它同时保住了节奏和错误，还关闭了对方主动解释的可能。\hciSee{2.9} 误区三。

\item 因为复述的作用是取得对方的确认，而不是展示你的理解。不复述、直接说“嗯”，对方不知道你懂没懂；复述完不等回答就继续，对方同样不知道，只是多了一层“他已经确认过”的错觉。两者在结果上等价。\hciSee{2.6}。

\item 因为“我没听懂”在上下级的通道里会被读成自我评价，而“我怕理解错”可以被读成责任心。两者传递的信息相同，后者成本更低。这不是虚伪，是在不允许示弱的通道里为必要的信息传输找到合法路径。\hciSee{2.8}。
\end{answerssec}

\hcirule

\begin{qabox}
\qaQ{读者提问}{我每次都当场发重传请求，但对方好像很不耐烦，怎么办？}
\qaA{本手册回答}{这是一个非常好的问题。需要说明的是，重传本身是有成本的，本手册的职责只是降低它的形式成本，不能消除它。若对方持续不耐烦，问题通常不在重传频率上，而在这段关系里“示弱”是否被允许——这超出了本手册的范围。具体效果因人而异。}
\end{qabox}

\hcirule

希望本章的内容能对你有所帮助。如需进一步了解第 3 章《确认》的相关内容，我可以随时为你展开。

%% 章末「页脚交互条」由 hci-manual.cls 自动挂接，无需手写。
